% !TEX root = ../esempio-documento.tex
% (la riga qui sopra dice a VS Code quale file compilare quando salvate questo)
% Capitolo richiamato da esempio-documento.tex con % !TEX root = ../esempio-documento.tex
% (la riga qui sopra dice a VS Code quale file compilare quando salvate questo)
% Capitolo richiamato da esempio-documento.tex con % !TEX root = ../esempio-documento.tex
% (la riga qui sopra dice a VS Code quale file compilare quando salvate questo)
% Capitolo richiamato da esempio-documento.tex con % !TEX root = ../esempio-documento.tex
% (la riga qui sopra dice a VS Code quale file compilare quando salvate questo)
% Capitolo richiamato da esempio-documento.tex con \input{capitoli/processi}

\section{Processi di supporto}

\subsection{Documentazione}
Ogni documento ha un redattore e un verificatore diversi. Il ciclo di
vita di un documento è:
\begin{enumerate}
  \item stesura nel proprio branch;
  \item apertura della pull request;
  \item verifica sul PDF dell'anteprima automatica;
  \item approvazione e unione in \glo{main}.
\end{enumerate}

\begin{nota}[Da ricordare]
La versione del documento è quella della prima riga del registro delle
modifiche: non va scritta da nessun'altra parte.
\end{nota}

\subsection{Versioni}
La Tabella~\ref{tab:versioni} riassume quando cambia ogni numero della
versione.

\begin{table}[H]
  \centering
  \caption{Regole per il numero di versione\label{tab:versioni}}
  \begin{tabular}{@{}ll@{}}
    \toprule
    \intestazione{Quando} & \intestazione{Nuova versione} \\
    \midrule
    Prima stesura                & 0.1.0 \\
    Ogni modifica prima della verifica & 0.2.0, 0.3.0, \ldots \\
    Verifica e approvazione      & 1.0.0 \\
    Modifica dopo l'approvazione & 1.1.0, 1.2.0, \ldots \\
    \bottomrule
  \end{tabular}
\end{table}

\subsection{Architettura di riferimento}
Le immagini stanno nella cartella \emph{allegati} e si inseriscono con
\emph{figura}, come in Figura~\ref{fig:architettura}.

\figura[0.9]{esempi/2026-10-15/architettura-aura.png}{Un'immagine di
esempio presa dalla cartella allegati\label{fig:architettura}}

\section{Verifica}
\subsection{Controlli sul documento}
Il Verificatore legge il PDF per intero e controlla ortografia, date,
nomi e riferimenti, poi approva la pull request.


\section{Processi di supporto}

\subsection{Documentazione}
Ogni documento ha un redattore e un verificatore diversi. Il ciclo di
vita di un documento è:
\begin{enumerate}
  \item stesura nel proprio branch;
  \item apertura della pull request;
  \item verifica sul PDF dell'anteprima automatica;
  \item approvazione e unione in \glo{main}.
\end{enumerate}

\begin{nota}[Da ricordare]
La versione del documento è quella della prima riga del registro delle
modifiche: non va scritta da nessun'altra parte.
\end{nota}

\subsection{Versioni}
La Tabella~\ref{tab:versioni} riassume quando cambia ogni numero della
versione.

\begin{table}[H]
  \centering
  \caption{Regole per il numero di versione\label{tab:versioni}}
  \begin{tabular}{@{}ll@{}}
    \toprule
    \intestazione{Quando} & \intestazione{Nuova versione} \\
    \midrule
    Prima stesura                & 0.1.0 \\
    Ogni modifica prima della verifica & 0.2.0, 0.3.0, \ldots \\
    Verifica e approvazione      & 1.0.0 \\
    Modifica dopo l'approvazione & 1.1.0, 1.2.0, \ldots \\
    \bottomrule
  \end{tabular}
\end{table}

\subsection{Architettura di riferimento}
Le immagini stanno nella cartella \emph{allegati} e si inseriscono con
\emph{figura}, come in Figura~\ref{fig:architettura}.

\figura[0.9]{esempi/2026-10-15/architettura-aura.png}{Un'immagine di
esempio presa dalla cartella allegati\label{fig:architettura}}

\section{Verifica}
\subsection{Controlli sul documento}
Il Verificatore legge il PDF per intero e controlla ortografia, date,
nomi e riferimenti, poi approva la pull request.


\section{Processi di supporto}

\subsection{Documentazione}
Ogni documento ha un redattore e un verificatore diversi. Il ciclo di
vita di un documento è:
\begin{enumerate}
  \item stesura nel proprio branch;
  \item apertura della pull request;
  \item verifica sul PDF dell'anteprima automatica;
  \item approvazione e unione in \glo{main}.
\end{enumerate}

\begin{nota}[Da ricordare]
La versione del documento è quella della prima riga del registro delle
modifiche: non va scritta da nessun'altra parte.
\end{nota}

\subsection{Versioni}
La Tabella~\ref{tab:versioni} riassume quando cambia ogni numero della
versione.

\begin{table}[H]
  \centering
  \caption{Regole per il numero di versione\label{tab:versioni}}
  \begin{tabular}{@{}ll@{}}
    \toprule
    \intestazione{Quando} & \intestazione{Nuova versione} \\
    \midrule
    Prima stesura                & 0.1.0 \\
    Ogni modifica prima della verifica & 0.2.0, 0.3.0, \ldots \\
    Verifica e approvazione      & 1.0.0 \\
    Modifica dopo l'approvazione & 1.1.0, 1.2.0, \ldots \\
    \bottomrule
  \end{tabular}
\end{table}

\subsection{Architettura di riferimento}
Le immagini stanno nella cartella \emph{allegati} e si inseriscono con
\emph{figura}, come in Figura~\ref{fig:architettura}.

\figura[0.9]{esempi/2026-10-15/architettura-aura.png}{Un'immagine di
esempio presa dalla cartella allegati\label{fig:architettura}}

\section{Verifica}
\subsection{Controlli sul documento}
Il Verificatore legge il PDF per intero e controlla ortografia, date,
nomi e riferimenti, poi approva la pull request.


\section{Processi di supporto}

\subsection{Documentazione}
Ogni documento ha un redattore e un verificatore diversi. Il ciclo di
vita di un documento è:
\begin{enumerate}
  \item stesura nel proprio branch;
  \item apertura della pull request;
  \item verifica sul PDF dell'anteprima automatica;
  \item approvazione e unione in \glo{main}.
\end{enumerate}

\begin{nota}[Da ricordare]
La versione del documento è quella della prima riga del registro delle
modifiche: non va scritta da nessun'altra parte.
\end{nota}

\subsection{Versioni}
La Tabella~\ref{tab:versioni} riassume quando cambia ogni numero della
versione.

\begin{table}[H]
  \centering
  \caption{Regole per il numero di versione\label{tab:versioni}}
  \begin{tabular}{@{}ll@{}}
    \toprule
    \intestazione{Quando} & \intestazione{Nuova versione} \\
    \midrule
    Prima stesura                & 0.1.0 \\
    Ogni modifica prima della verifica & 0.2.0, 0.3.0, \ldots \\
    Verifica e approvazione      & 1.0.0 \\
    Modifica dopo l'approvazione & 1.1.0, 1.2.0, \ldots \\
    \bottomrule
  \end{tabular}
\end{table}

\subsection{Architettura di riferimento}
Le immagini stanno nella cartella \emph{allegati} e si inseriscono con
\emph{figura}, come in Figura~\ref{fig:architettura}.

\figura[0.9]{esempi/2026-10-15/architettura-aura.png}{Un'immagine di
esempio presa dalla cartella allegati\label{fig:architettura}}

\section{Verifica}
\subsection{Controlli sul documento}
Il Verificatore legge il PDF per intero e controlla ortografia, date,
nomi e riferimenti, poi approva la pull request.
